\documentclass[11pt]{article}

\usepackage[margin=25mm]{geometry}

\usepackage{amsmath,amssymb,fontspec,unicode-math}

\IfFontExistsTF{DejaVu Serif}{\setmainfont{DejaVu Serif}}{\setmainfont{Latin Modern Roman}}

\IfFontExistsTF{Latin Modern Math}{\setmathfont{Latin Modern Math}}{\setmathfont{DejaVu Math TeX Gyre}}

\usepackage{xurl,hyperref}

\setlength{\parindent}{0pt}\setlength{\parskip}{6pt}

\begin{document}

\title{Visual Jenga: Discovering Object Dependencies via Counterfactual Inpainting}

\author{ELVIA preliminary AI analysis}\date{21 October 2026}\maketitle

Preliminary assessment: No Issue Found.\par

Independent scientific review remains pending. Author adoption has not been requested.

No substantive implementation or empirical contradiction was established. The conditional-probability comparison and feature score have mathematical scope limits consistent with the paper’s explicit co-occurrence-based heuristic and lack of causal reasoning.

\section*{Review scope}

Claim-focused reading of the original paper and relevant author-released passages; direct mathematical or static-source reasoning. No author program, earlier verifier, model or benchmark was executed during this check.

\section*{Problem 1: Conditional asymmetry does not identify a physical cause}

Source: Section 3.1, Eq.(1), PDF p.4; Limitations, PDF p.9

Assessment: not\_issue; impact: minor

\subsection*{Source claim (paraphrase)}

The asymmetry \(P(A\mid B)\ll P(B\mid A)\) motivates a directional dependency heuristic.

\subsection*{Analysis}

For positive joint probability, \(P(A\mid B)/P(B\mid A)=P(A)/P(B)\). Independent events with marginals \(0.01\) and \(0.9\) have joint \(0.009\), giving the large conditional asymmetry despite independence. The same identity holds under a common fixed context. The curated benchmark additionally selects physical-support pairs, and the paper explicitly describes its approach as statistical co-occurrence without causal reasoning.

\subsection*{Finding and scope}

The comparison alone is not causal identification. This scoped limitation does not contradict the qualified heuristic or its reported benchmark arithmetic.

\section*{Problem 2: Mean similarity does not uniquely specify distributional peakedness}

Source: Section 3.3(ii), Eq.(2), PDF p.5

Assessment: not\_issue; impact: minor

\subsection*{Source claim (paraphrase)}

The product of mean CLIP/DINO similarities is motivated by peakedness of the inpainting distribution.

\subsection*{Analysis}

In both feature spaces take original unit vector \((1,0)\) and candidate outputs \((0.6,0.8)\), \((0.6,-0.8)\), with area fraction \(1\). Both rescaled similarities are \(0.8\). A point mass and an equal two-point mixture have entropy \(0\) and \(\log2\), but score \(1-0.8^2=0.36\) in both cases. Real-image realizability of these simultaneous embeddings is unestablished.

\subsection*{Finding and scope}

The score measures deviation from the original crop and need not identify entropy. Its empirical utility is not contradicted.

\section*{Limits}

Preliminary AI assessment; independent human scientific adjudication remains pending.\par

Analytic witnesses and static release checks do not remeasure the published benchmark results.\par

\section*{Source identity}

NeurIPS 2025 official proceedings PDF (retained local copy)

\url{https://proceedings.neurips.cc/paper_files/paper/2025/file/ea7efe35cef8808c2fa66c349671a73d-Paper-Conference.pdf}

Reviewed PDF SHA-256: \nolinkurl{114498f81739dd19ef17a8a82f229d02fec139671bf38945304499163bdfe9ae}

The SHA-256 identifies the exact retained PDF inspected in this rapid source/argument check. Original supplied report bytes are retained locally. The original PDF retrieval timestamp is taken from matching retained metadata.

\end{document}
